\begin{remark}[Scheme dependence]
\label{rem:c10-scheme1}
(a) $\hbar F_1^{(1)}(x)$ enters the one-body kernel alone, through $-\{\hbar F_1^{(1)}(x),-\}$. (b) $\hbar F_1^{(0)}$ enters through trees, one-body. (c) A finite part $\psi'$ of the operators changes $L_1$ by $[s,\psi']$ (Proposition \ref{prop:c8-laws} (b)), and the component along $x$ by $[\{S^{(1)}(x),-\},\psi']$; for $\psi'=\tfrac a2\{S^{(1)}(\Delta),-\}$ by $\tfrac a2\lambda_x\{S^{(1)}(x),-\}$ (Proposition \ref{prop:c9-scheme}). (d) The cutoff shape enters the finite parts of the special-edge cycles (Proposition \ref{prop:c10-shape}, Lemma \ref{lem:c10-special}). The finite vertex \eqref{eq:c10-finitesp}, the coefficient of $\varepsilon^0$ of the special-edge parts of the counterterm brackets, is cancelled in the components by the graphs with the special edge on a bridge: a bridge and its counterterm partner give $e^{-\varepsilon|p|^2}-1$ times a renormalised subgraph, which tends to $0$ (Theorem \ref{thm:c11-classes} (ii)); the two-body kernels are the coefficients of $\varepsilon^0$ of the defect cycles alone, in arity one and in arity zero. The spectrum of $D_1$ on cohomology, at the coefficients of $\hbar F_1^{(1)}(\Delta)$ fixed by the arity-two equations (Theorem \ref{thm:c9-theta2}), is independent of (a)--(d) (Theorem \ref{thm:c8-invariants} (iii)).
\end{remark}
